\documentclass[12pt, letterpaper]{article}
\usepackage{amsthm, amsmath, amssymb, graphicx, hyperref, url, pdfpages, setspace}
\onehalfspacing
\usepackage[utf8]{inputenc}
\usepackage{pgfplots}
\pgfplotsset{compat=1.18}
\usepackage{tikz}
\usetikzlibrary{arrows.meta, positioning}

\newtheorem{thm}{Theorem}[section]
\newtheorem{lemma}[thm]{Lemma}

\theoremstyle{definition}
\newtheorem{ex}[thm]{Example}
\newtheorem{defn}[thm]{Definition}

\title{Power Series}
\author{Gordon Chen}
\date{}

\begin{document}
\maketitle

\section{Weierstrass M-Test}
\begin{thm}
  Let $f_j : X \to \mathbb{R}$.
  If $|f_j| \le M_j$ on $X$ and $\sum_{j=0}^\infty M_j$ converges,
  then the series $\sum_{j=0}^\infty f_j(x)$ converges absolutely and uniformly on $X$.

  \begin{proof}
    Let $\epsilon > 0$. Since $\sum_{j=0}^\infty M_j$ converges, there exists some $J$ such that
    $n > m \ge J$ implies that $\sum_{j=m+1}^n M_j < \epsilon$. Then 
    $$
    |s_n(x) - s_m(x)| = \left|\sum_{j=m+1}^n f_j(x) \right|
    \le \sum_{j=m+1}^n |f_j(x)| \le \sum_{j=m+1}^n M_j
    $$
    for $n > m \ge J$ and all $x \in X$, so the sequence of partial sums is uniformly Cauchy
    on complete $\mathbb{R}$, so the series converges uniformly
    (and absolutely since we bounded $\sum_{j=m+1}^n |f_j(x)|$).
  \end{proof}
\end{thm}


\section{Power Series}
\begin{thm}
  The power series $\sum_{j=0}^\infty a_j (x-b)^j$ converges absolutely on
  the interval $(b-R, b+R)$ where $R = 1 / \limsup_j |a_j|^{1/j}$ is
  the radius of convergence.
  Furthermore, the power series converges uniformly on and is continuous
  on every compact subinterval $[b-r, b+r]$, where $0 \le r < R$.

  \begin{proof}
    By the Root Test, the power series $\sum_{j=0}^\infty a_j (x-b)^j$ converges when
    $\limsup_j |a_j (x-b)^j|^{1/j} < 1,$
    or equivalently, when $|x-b| < 1 / \limsup_j |a_j|^{1/j}$.
    We denote the radius of convergence $R = 1 / \limsup_j |a_j|^{1/j}$.
    Then, by the Root Test, the power series converges on $(b-R, b+R)$ and diverges outside
    $[b-R, b+R]$.

    Let $|x-b| \le r < R$.
    Then $|a_j (x-b)^j| \le |a_j|r^j = M_j$.
    Applying the Root Test, we get that
    $\limsup_j M_j^{1/j} = \limsup_j |a_j|^{1/j}r = \frac{r}{R} < 1$,
    so $\sum_{j=0}^\infty M_j$ converges.
    Then, by the Weierstrass M-Test, we have that $\sum_{j=0}^\infty a_j (x-b)^j$ converges
    uniformly on every compact subinterval $[b-r, b+r]$, where $0 \le r < R$.

    Since the series converges uniformly on every compact subinterval
    and the series is the limit of partial sums which are polynomials and therefore
    continuous functions, then the series is continuous on any compact subinterval.
    Furthermore, since every point $x \in (b-R, b+R)$ is in
    some compact subinterval, the series is continuous on $(b-R, b+R)$.
  \end{proof}
\end{thm}


\section{Differentiating and Integrating}
\begin{thm}
  The power series $\sum_{j=0}^\infty a_j (x-b)^j$ converging on $(b-R, b+R)$ can be
  integrated and differentiated term-by-term infinitely many times on any compact
  subinterval of $(b-R, b+R)$.

  \begin{proof}
    Consider the series we get by term-by-term integration and differentiation,
    $\sum_{j=0}^\infty \frac{a_j}{j+1} (x-b)^{j+1}$ and $\sum_{j=1}^\infty j a_j (x-b)^{j-1}$
    respectively.
    Using the fact that $\lim_j \left(\frac{1}{j+1}\right)^{1/j} = \lim_j j^{1/j} =1$,
    applying the Root Test to the new series gives
    $$
      \limsup_j \left|\frac{a_j}{j+1} (x-b)^{j+1}\right|^{1/j}
      = \limsup_j \left(\frac{1}{j+1}\right)^{1/j} |a_j|^{1/j} |x-b|^{1+1/j}
      = \frac{|x-b|}{R}
    $$
    and
    $$
      \limsup_j \left|j a_j (x-b)^{j-1}\right|^{1/j}
      = \limsup_j j^{1/j} |a_j|^{1/j} |x-b|^{1-1/j}
      = \frac{|x-b|}{R}
    $$
    so the new series converge on $(b-R, b+R)$ and diverge outside $[b-R, b+R]$
    just like the original series,
    so they are also uniformly convergent on any compact subinterval of $(b-R, b+R)$.

    Uniform convergence of the series then allows us to integrate term by term and get the integral
    of the entire series (where the integral is evaluated between any two points inside an interval
    of uniform convergence):
    $$
    \int_b^x \sum_{j=0}^\infty a_j (t-b)^j dt = \sum_{j=0}^\infty \int_b^x a_j (t-b)^j dt
    = \sum_{j=0}^\infty \frac{a_j}{j+1} (x-b)^{j+1}.
    $$

    Since all partial sums are differentiable (polynomials),
    the derivatives of the partial sums converge uniformly on every compact subinterval,
    and the series is convergent at at least one point,
    then the series is differentiable and its derivative (at a point inside the subinterval)
    is the term by term derivative (the limit of the derivatives of the partial sums):
    $$
    \frac{d}{dx} \sum_{j=0}^\infty a_j (x-b)^j = \sum_{j=0}^\infty \frac{d}{dx} a_j (x-b)^j
    = \sum_{j=1}^\infty j a_j (x-b)^{j-1}.
    $$

    Since the radius of convergence of the term-by-term differentiated and integrated series
    is the same as the original series, we can differentiate and integrate again, term-by-term, and
    infinitely many times on any compact interval inside $(b-R, b+R)$.
  \end{proof}
\end{thm}


\section{Examples}
\begin{ex}
  Let $e^x$ be defined as $e^x = \sum_{j=0}^\infty \frac{x^j}{j!}$.

  By the Root Test, $e^x$ converges when
  $\limsup_j \left|\frac{x^j}{j!}\right|^{1/j} = |x| \limsup_j \frac{1}{(j!)^{1/j}} < 1$.
  See that $j! \ge (\frac{j}{2})^\frac{j}{2}$, so
  $(j!)^{1/j} \ge (\frac{j}{2})^{\frac{1}{2}}$, and $\lim_j (j!)^{1/j} = \infty$.
  Then $\lim_j \frac{1}{(j!)^{1/j}} = 0$, so $e^x$ converges on $\mathbb{R}$ by the
  Root Test.

  Then $e^x$ converges uniformly on every compact subinterval of $\mathbb{R}$,
  so we can differentiate and integrate it term-by-term infinitely many times on that subinterval.
  Then $e^x$ is infinitely differentiable on $\mathbb{R}$ and integrable on every
  compact subinterval of $\mathbb{R}$.
\end{ex}

\end{document}
